\begin{table*}[t]
\centering
\small
\caption{One-shot false-feasible traversal success does not establish correct rejection. This table separates explicit abstention from collision and non-collision execution failure.}
\label{tab:false-feasible-outcomes}
\begin{tabular}{lrrrrrrrr}
\toprule
Method & Episodes & Traversal success & Explicit reject & Correct reject & Collision & Near collision & Timeout/stuck & Wasted attempts \\
\midrule
Reactive rule baseline & 1500 & 0.0$\pm$0.0\% (0/1500) & 0.0$\pm$0.0\% (0/1500) & 0.0$\pm$0.0\% (0/1500) & 100.0$\pm$0.0\% (1500/1500) & 57.6$\pm$2.6\% (864/1500) & 0.0$\pm$0.0\% (0/1500) & 100.0$\pm$0.0\% (1500/1500) \\
DEGNAV-Rule & 1500 & 0.0$\pm$0.0\% (0/1500) & 0.0$\pm$0.0\% (0/1500) & 0.0$\pm$0.0\% (0/1500) & 2.8$\pm$0.7\% (42/1500) & 14.5$\pm$1.0\% (218/1500) & 97.2$\pm$0.7\% (1458/1500) & 100.0$\pm$0.0\% (1500/1500) \\
\bottomrule
\end{tabular}
\vspace{0.25em}
\begin{minipage}{0.98\linewidth}
\footnotesize Explicit reject means the controller selected Reject; correct reject is explicit rejection on a benchmark-labeled infeasible passage. Timeout/stuck is not safe rejection. Collision and near-collision are logged independently and may overlap. Recurrence handling is evaluated separately through memory.
\end{minipage}
\end{table*}
